\subsubsection{Why Teleportation Works}\label{sec:why-teleportation-works}

Now that we've completed the mathematics, this last section should answer a single question: \textbf{\emph{why didn't we break any law of quantum mechanics?}} Let's revisit each principle.

\highspace
\begin{flushleft}
    \textcolor{Green3}{\faIcon{check-circle} \textbf{The quantum state is transferred, not copied}}
\end{flushleft}
We could think that we copied the state of Alice's qubit to Bob's qubit:
\begin{equation*}
    \text{Alice: } \ket{\psi} \quad \text{Bob: } \ket{0} \quad \xrightarrow{\text{teleportation}} \quad \text{Alice: } \ket{\psi} \quad \text{Bob: } \ket{\psi}
\end{equation*}
But we didn't. Teleportation does \textbf{not} do this, instead it destroys the state of Alice's qubit and transfers it to Bob's qubit:
\begin{equation*}
    \text{Alice: } \ket{\psi} \quad \text{Bob: Bell qubit} \quad \xrightarrow{\text{teleportation}} \quad \text{Alice: (measured)} \quad \text{Bob: } \ket{\psi}
\end{equation*}
Alice no longer possesses the state $\ket{\psi}$, and Bob does. The state has been \textbf{moved}, not duplicated. This is why teleportation \textbf{does not violate the no-cloning theorem}.

\highspace
\begin{flushleft}
    \textcolor{Green3}{\faIcon{question-circle} \textbf{Why are the two classical bits necessary?}}
\end{flushleft}
Suppose Alice refuses to send the message to Bob. Then Bob waits forever because he only knows that he has a Bell qubit, which is one of the four possible states: \ket{\psi}, $X\ket{\psi}$, $Z\ket{\psi}$, or $XZ\ket{\psi}$. But he has no idea which one. He cannot guess, nor can he measure his qubit to find out, because measuring it would destroy the state. He is stuck.

\highspace
Only when Alice sends 00, 01, 10, or 11 can Bob apply the correct Pauli gate to recover the state $\ket{\psi}$. This is why \textbf{the two classical bits are necessary}. They do not encode the amplitudes $\alpha$ and $\beta$, but they encode the \textbf{correction} that Bob must apply to his qubit to recover the state $\ket{\psi}$.

\highspace
\begin{flushleft}
    \textcolor{Green3}{\faIcon{link} \textbf{Why doesn't this allow faster-than-light communication?}}
\end{flushleft}
At first glance it looks like something magical happened: Alice measures; instantly Bob's qubit becomes 1 of the 4 states. \emph{\textbf{Could Alice send information faster than light?}} The answer is obviously no.

\highspace
Imagine Alice measures her qubit. Bob immediately looks at his own qubit. Can he tell whether Alice obtained 00 or 11? Or 01 or 10? No, he cannot. Can he determine whether he should apply gate $X$, $Y$ or $Z$, or nothing? No, he cannot. Thus, can he recover the original state \ket{\psi} without Alice? No! He must wait for the classical message from Alice. This is why \textbf{teleportation does not allow faster-than-light communication}. Classical communication is still necessary to complete the teleportation process, and it cannot travel faster than light.

\highspace
\begin{flushleft}
    \textcolor{Red2}{\faIcon{exclamation-triangle} \textbf{The Bell pair is a unique resource}}
\end{flushleft}
At the beginning we had the resource:
\begin{equation*}
    \ket{\Phi^+} = \frac{1}{\sqrt{2}}\left(\ket{00} + \ket{11}\right)
\end{equation*}
One might ask: ``\emph{can Alice and Bob teleport another qubit using the same Bell pair?}''. The answer is \textbf{no}. After Alice's measurement, the entanglement is gone. The Bell pair has been ``\emph{spent}'' and cannot be reused. If Alice and Bob want to teleport another qubit, they must first create a \textbf{new Bell pair}. It is a sort of battery: before teleportation it is fully charged (shared entanglement); during teleportation its energy is used; after teleportation it is discharged and cannot be reused. This is why \textbf{the Bell pair is a unique resource}.

\highspace
\begin{flushleft}
    \textcolor{Green3}{\faIcon{question-circle} \textbf{If we still have to send classical information, what's the point of teleportation?}}
\end{flushleft}
Let's imagine the classical world. Suppose Alice wants Bob to have \ket{\psi}. Could she simply send him the values $\alpha$ and $\beta$ over the internet? No, because \textbf{she doesn't know them}. Remember that Alice never learns $\alpha$, $\beta$, because if she did, the state would collapse and she would lose the information.

\highspace
Now, let's think about this more practically. Suppose Alice has a quantum state:
\begin{equation*}
    \ket{\psi} = 0.834\ket{0} + 0.551e^{i 0.42}\ket{1}
\end{equation*}
\emph{How would Alice tell Bob this?} She would have to send $\alpha = 0.834$ and $\beta = 0.551e^{i 0.42}$, but she cannot because the qubit is unknown to her. She has only one copy of the qubit, so measuring destroys it. So classical communication alone is useless.

\highspace
The advantage is \textbf{not faster communication}. It is something much stranger: quantum mechanics lets us transfer an \textbf{unknown quantum state} from one qubit to another, without ever knowing what the state is. This is something that is impossible in classical physics.

\highspace
\textcolor{Green3}{\faIcon{question-circle} \textbf{Where is this useful?}} If teleportation were only a way of sending qubits from one room to another, it would be an expensive curiosity. Its real applications are in \textbf{quantum networks}.

\highspace
Imagine a situation where there are two quantum computers, one in Milan (Quantum Computer A) and one in New York (Quantum Computer B). They share a Bell pair. Suppose computer A has produced a qubit during some computation. Instead of physically transporting the fragile qubit through a noisy optical fiber, it can teleport its state to quantum computer B. This is one of the building blocks of the future \textbf{quantum internet}.

\highspace
How two quantum computers can share a Bell pair is a topic beyond the scope of this course. In practice, one half of an entangled Bell pair can be physically distributed to a remote quantum computer using quantum communication channels, such as optical fibers or satellite links. Establishing and maintaining long-distance entanglement is one of the central challenges in the development of the future quantum internet and remains an active area of research.